\section{The numerical comparisons and the bound 498}\label{sec:completion}
We now assemble the proof. There are two parts to the numerical argument. A uniform estimate, supplemented by adjacent
interval comparisons, excludes every $D\geq875$. Below $875$ the elementary
reduction leaves sixteen possible values, each of which is excluded by a
separate determinant comparison. We describe the coverage first, then give
the parameters and the decisive calculation at $499$.

\subsection{The range \texorpdfstring{$D\geq875$}{D >= 875}}
For this range we use the resultant rank criterion of
Proposition~\ref{prop:resultant-rank} and the radial size estimate of
Proposition~\ref{prop:radial-size}. Appendix~\ref{app:global} applies these
estimates on an unbounded interval and then on finitely many adjoining
intervals below it. Set $D_0=5\,760\,000$, $X=(D-3)/4$, and
$X_0=(D_0-3)/4$. On the unbounded interval, it gives
\begin{equation}\label{eq:tail-gap}
 G_{\rm tail}(D)
 >0.00031743+0.127829\log(X/X_0)\qquad(D\geq D_0),
\end{equation}
where $G_{\rm tail}$ is the normalized difference between the arithmetic
lower bound and the radial upper bound. This difference is positive
throughout $D\geq D_0$, contradicting the determinant inequality.

The interval between $875$ and $D_0$ is covered as follows. The gap bounds in
Table~\ref{tab:coverage} are rounded downward from the smallest
lower endpoint in each family. Each comparison uses a closed interval, and
successive intervals share an endpoint.
\begin{table}[ht]
\centering
\begin{tabular}{lrr}
\toprule
Range & Number of intervals & Lower gap bound\\
\midrule
$[875,3\,968\,539]$ & 28 & $0.00397338$\\
$[3\,968\,539,4\,894\,699]$ & 37 & $0.00032629$\\
$[4\,894\,699,5\,760\,000]$ & 36 & $0.00040940$\\
\bottomrule
\end{tabular}
\caption{Interval comparisons covering $[875,D_0]$. The first
family uses norm-form estimates, the second exact finite prime sums, and the
third explicit prime-sum bounds of Rosser and Schoenfeld.}
\label{tab:coverage}
\end{table}

Thus \eqref{eq:tail-gap} and
Table~\ref{tab:coverage} exclude all $D\geq875$ under the class-number-one
hypothesis.

\subsection{Parameters and rank for the remaining comparisons}
By Lemma~\ref{lem:reduction}, the values still possible above $498$ are $499$
and the fifteen primes in Table~\ref{tab:candidates}. For these comparisons
we use $K=43$ and a prime cutoff $P$. The other parameters are
\[
 S=\left\lfloor\frac{6P}{5K}\right\rfloor,
 \qquad R=KS,\qquad J=\lfloor R/8\rfloor,
 \qquad M=R(S+1)/2.
\]
For $D=499$ we take $P=312\,000\,000$ and $n=35\,394\,400$.
Table~\ref{tab:parameters} gives all four parameter choices.

\begin{table}[ht]
\centering
\begin{tabular}{rrrrr}
\toprule
$P$ & $S$ & $R$ & $J$ & $n$\\
\midrule
$36\,000\,000$ & $1\,004\,651$ & $43\,199\,993$ & $5\,399\,999$ & $4\,822\,400$\\
$169\,000\,000$ & $4\,716\,279$ & $202\,799\,997$ & $25\,349\,999$ & $22\,638\,400$\\
$312\,000\,000$ & $8\,706\,976$ & $374\,399\,968$ & $46\,799\,996$ & $41\,793\,600$\\
$312\,000\,000$ & $8\,706\,976$ & $374\,399\,968$ & $46\,799\,996$ & $35\,394\,400$\\
\bottomrule
\end{tabular}
\caption{Parameters for all determinant comparisons below $875$. The last
line is used on $[499,503]$; the other lines at the individual discriminants
listed in Table~\ref{tab:candidates}.}
\label{tab:parameters}
\end{table}

We first prove nonvanishing. There are $814$ primes in
$(A,B]=(160\,000,170\,000]$. Lemma~\ref{lem:auxiliary-count} bounds the number of non-inert primes
by a sum of square roots. This sum is less than $765$, both on $[499,503]$
and at each value in Table~\ref{tab:candidates}. Hence at least one of the
$814$ primes in the interval is inert.

Lemma~\ref{lem:node-disk} gives $2\rho_n<28\,928$ for every comparison
in Table~\ref{tab:parameters}, and $2\rho_n<22\,492$ at $D=499$.
Both are below $A=160\,000$, so Corollary~\ref{cor:distinct-residues}
gives distinct residues modulo every prime in $(A,B]$. Moreover,
$K<A<B<R<A^2$ in every case. The remaining rank conditions are
$J\leq R$, $n\geq2S$, and
\begin{equation}\label{eq:uniform-rank-check}
 n\left(J-\frac{RJ}{A^2}-\frac RA\right)-\frac{(R+B)^2}{2K}>0.
\end{equation}
At the smaller node count the difference in \eqref{eq:uniform-rank-check}
exceeds $7.22\times10^{11}$.
The reason for narrowing the prime interval can now be seen quantitatively.
Writing $\theta=M/(nJ)$, inequality~\eqref{eq:uniform-rank-check} is equivalent to
\[
 \theta<\frac{(1+1/S)(1-R/A^2-R/(AJ))}{(1+B/R)^2}.
\]
At the parameters for $499$, the right side is approximately $0.984431$
when $A=160\,000$, just above the actual occupancy $0.983995$. With
$A=150\,000$ it is approximately $0.982414$, which is too small.
The narrower interval improves the sufficient rank estimate while still
containing more primes than the norm form can account for as non-inert.

At the other three parameter choices the respective differences exceed
$4.12\times10^{12}$, $9.02\times10^{13}$ and $2.95\times10^{14}$.
Theorem~\ref{thm:simple-rank} therefore proves full row rank in every case.

\subsection{The fifteen comparisons above 499}
We now evaluate Proposition~\ref{prop:comparison}, with the local arithmetic
estimates of Proposition~\ref{prop:general-arithmetic} and the spectral bound
\eqref{eq:metric-mesh}. The spectral parameters and mesh sizes are fixed as
in \eqref{eq:spectral-parameters}. Table~\ref{tab:candidates} gives the result;
each displayed lower bound is rounded downward. The calculation uses outward rounding.
\begin{table}[ht]
\centering
\begin{tabular}{rrr@{\qquad}rrr}
\toprule
$D$ & $P/10^6$ & Lower gap bound & $D$ & $P/10^6$ & Lower gap bound\\
\midrule
523&312&0.01913120 &683&36&0.24689369\\
547&312&0.06770128 &691&36&0.25866104\\
563&169&0.09816857 &739&36&0.33084083\\
571&169&0.10589325 &787&36&0.39774315\\
587&36&0.08275416 &811&36&0.42960657\\
619&36&0.13882635 &827&36&0.45047644\\
643&36&0.17968979 &859&36&0.48970182\\
659&36&0.20911703 &&&\\
\bottomrule
\end{tabular}
\caption{Positive lower bounds for $G=\arith-\metric$ at all fifteen primes
$D\equiv3\pmod8$ in $(499,875)$.}
\label{tab:candidates}
\end{table}
Each entry is positive and contradicts \eqref{eq:master}. Thus every value in
the table is excluded.

\subsection{The decisive comparison at 499}
At $D=499$ we retain $P=312\,000\,000$, $K=43$ and the same $R,S,J$ as in
the preceding comparison, but reduce the node count to $n=35\,394\,400$.
The arithmetic and size calculations are performed on the whole interval
$[499,503]$, using the same parameters throughout. Only $499$ needs to be excluded: $503\equiv7\pmod8$ is already
impossible by Lemma~\ref{lem:reduction}.

Table~\ref{tab:499} separates the three contributions to the comparison.
The common term is
\[
 U_{\rm common}=\frac{W_1+W_2}{W}\log\rho+
                     \frac{Mc+2M\log M}{W},
\]
so that $G=\arith-U_{\rm common}-U_{\rm var}$.
The displayed arithmetic numbers are lower bounds and the size numbers are
upper bounds, rounded in the safe direction. The first numerical column uses the larger node count, with all
other parameters fixed.
\begin{table}[ht]
\centering
\begin{tabular}{lrr}
\toprule
 & Larger node count & Smaller node count\\
\midrule
$n$ & $41\,793\,600$ & $35\,394\,400$\\
$\theta=M/(nJ)$, approximately & $0.83333113$ & $0.98399487$\\
Arithmetic lower bound $\arith$ & $9.99153787$ & $9.99493844$\\
Upper bound for $U_{\rm common}$ & $9.71143220$ & $9.62002864$\\
Upper bound for $U_{\rm var}$ & $0.40353803$ & $0.32371122$\\
\bottomrule
\end{tabular}
\caption{The comparison on $[499,503]$ with $M=1\,629\,945\,955\,088\,368$.
Only the occupancy row uses approximate decimals; the other decimal entries
are one-sided bounds.}
\label{tab:499}
\end{table}

The arithmetic decomposition \eqref{eq:arithmetic-parts} at the smaller
node count is, approximately,
\[
 \arith\approx15.49500-6.91595+1.37599+0.03990\approx9.99494.
\]
The last two contributions are the two- and three-factor corrections.
Together they restore about $1.416$ to the arithmetic lower bound, far more
than the final margin; even the three-factor term, about $0.040$, is
comparable to that margin.

The change in the gap is about $0.1746$. Of this, $0.0914$ comes from the
common term, $0.0798$ from the variable spectral term, and $0.0034$ from the
arithmetic lower bound. The common-term improvement is close to the radius
estimate $0.085$ in Section~\ref{subsec:central}. The variable term benefits
from the higher occupancy and the changed residual scale $\ell$. 

The complete directed calculation gives the following enclosure for the
computed gap bound at the smaller node count, rounded outward here:
\begin{equation}\label{eq:499-gap}
 G\in[0.05119859,\;0.05120060].
\end{equation}
In particular $G>0.05119859$, contradicting the determinant comparison.
For the larger node count the corresponding enclosure is
\[
 [-0.12343235,\;-0.12343034].
\]

\begin{proof}[Proof of Theorem~\ref{thm:main}]
Suppose that $h(-D)=1$ and $D>498$. Lemma~\ref{lem:reduction} makes $D$ a
prime congruent to $3$ modulo $8$. Equation~\eqref{eq:tail-gap} and
Table~\ref{tab:coverage} exclude $D\geq875$. Table~\ref{tab:candidates}
excludes all the remaining possibilities greater than $499$, and
\eqref{eq:499-gap} excludes $499$ itself. These cases exhaust all $D>498$.
\end{proof}

At $499$, the occupancy is already about $0.984$, so the number of nodes
is only about $1.6\%$ above the dimension-count minimum $M/J$. There is
therefore little room to reduce the radius further by improving rank alone:
in the radius approximation of Section~\ref{subsec:central}, reaching the
dimension count exactly would save only about $0.008$ in the common term.
Sharper arithmetic or spectral estimates are natural candidates for further
progress.

The class number one hypothesis is used to turn non-inert prime ideals into
representations by the norm form. A comparable estimate for general
Landau--Siegel zeros would require a substitute for this arithmetic input.
